\chapter{Anàlisi Lèxica i Sintàctica}
L'anàlisi d'un programa font es divideix tradicionalment en dues fases clarament diferenciades: l'anàlisi lèxica (\textit{lexing}) i l'anàlisi sintàctica (\textit{parsing}). En el disseny d'aquest intèrpret, s'ha optat per defugir l'ús d'eines externes generadores de codi com \textit{Alex} o \textit{Happy}. En el seu lloc, s'ha implementat un enfocament natiu basat en \textbf{combinadors de parsers} encastats directament en el llenguatge amfitrió, Haskell. Aquesta decisió de disseny permet mantenir l'homogeneïtat absoluta de la base de codi, aprofita el sistema de tipus de Haskell per garantir la seguretat en temps de compilació i simplifica dràsticament l'arquitectura de l'intèrpret.

Per tal d'il·lustrar de manera didàctica i cohesiva el funcionament intern del sistema, aquesta memòria utilitzarà una \textbf{expressió guia} que ens acompanyarà de manera transversal al llarg de tots els capítols. Analitzarem el ``viatge'' complet d'aquesta línia, des que és una simple cadena de text pla fins a la seva reducció final en la màquina de grafs. L'expressió triada és:

\begin{center}
    \texttt{(\textbackslash x y -> x * y) 2 3}
\end{center}

En aquest capítol, s'exposa com aquesta seqüència de caràcters es transforma en una estructura de dades arbòria que reflecteix la semàntica formal del llenguatge.

\section{Disseny de la Gramàtica i l'Arbre de Sintaxi Abstracta}

L'objectiu primer d'aquesta fase és construir una representació abstracta en forma d'arbre (AST, \textit{Abstract Syntax Tree}) de qualsevol codi subministrat per l'usuari. El tipus de dades algebraic central que governa el nucli del llenguatge és \texttt{Expr}, definit al mòdul \texttt{HSTipus}:

\begin{lstlisting}[style=haskellstyle, caption=Definició del tipus algebraic d'expressions (Expr)]
data Expr = Val Int                 
          | Var String              
          | App Expr Expr           
          | Lam String Expr       
          | If Expr Expr Expr
          deriving (Show,Eq)
\end{lstlisting}

\subsection{Suficiència Teòrica del Nucli d'Expressions}
Aquest conjunt de cinc nodes és minimalista però completament suficient per formalitzar un llenguatge de programació funcional complet. El fonament teòric d'aquesta suficiència se sustenta en l'estratègia exposada per Simon Peyton Jones \cite{peyton1987implementation}, on es defensa que un llenguatge funcional d'alt nivell s'ha de traduir i reduir a un \textbf{Càlcul Lambda Enriquit} (\textit{Enriched Lambda Calculus}) abans de ser executat. 

El Càlcul Lambda pur només necessita tres constructors fonamentals: variables (\texttt{Var}), abstraccions (\texttt{Lam}) i aplicacions (\texttt{App}), els quals defineixen un sistema Turing complet capaç de modelar qualsevol tipus de dada mitjançant codificacions estructurals. No obstant això, tal com argumenta Peyton Jones \cite{peyton1987implementation}, per motius d'eficiència pràctica de l'intèrpret s'estén aquest nucli amb constants primitives (\texttt{Val Int}) i estructures de control de flux primitives (\texttt{If Expr Expr Expr}). Els operadors aritmètics i lògics (com \texttt{*} o \texttt{+}) no requereixen nodes especialitzats a l'AST; es tracten formalment com a variables (\texttt{Var}) el tipus funcional de les quals es troba predefinit a l'entorn global o \textit{Prelude}, d'acord amb el model de funcions primitives integrades descrit a \cite{peyton1987implementation}.

\subsection{El Tractament de la Recursivitat mitjançant Supercombinadors}
Un detall de disseny crucial heretat de les tesis de Peyton Jones \cite{peyton1987implementation} és l'absència d'un node de recursivitat dedicat dins de \texttt{Expr} (com podria ser un constructor \texttt{LetRec} local o un operador de punt fix \texttt{Fix}). 

En el nostre intèrpret, la recursivitat no es gestiona a nivell d'expressió pura, sinó a nivell global mitjançant la transformació del programa en \textbf{Supercombinadors}. Quan el parser detecta definicions de funcions recursives, les lliga a l'entorn de declaracions. Més endavant, el motor de l'intèrpret tradueix aquestes declaracions en supercombinadors globals on els cicles i l'auto-referència es resolen mitjançant la instanciació de punters de cicle directes en el \textit{heap}, evitant la necessitat d'introduir constructs complexos de fixació local a l'AST durant la inferència de tipus.

\section{L'Abstracció de Declaracions (\textit{Bindings})}

Un programa real no és només una expressió aïllada; es compon d'un conjunt de definicions globals, de vegades recursives, seguides de l'expressió que es vol avaluar. Per modelar aquesta estructura, el parser introdueix els conceptes de \texttt{Binding} i els combinadors de nivell superior \texttt{program} i \texttt{line}.

\begin{lstlisting}[style=haskellstyle, caption=Estructures d'enllaç i flux del programa]
data Binding = Bind String Expr deriving (Show)
type Program = ([Binding], Expr)
\end{lstlisting}

El combinador \texttt{binding} realitza el procés clàssic de \textbf{desugarejament} (\textit{desugaring}) fonamental per transformar el codi amigable d'estil Haskell cap al nucli del Càlcul Lambda Enriquit \cite{peyton1987implementation}. Permet dues estructures clau:

\begin{enumerate}
    \item \textbf{Definicions Estàndard Multiargument:} Permet escriure definicions lineals com \texttt{suma x y = x + y}. El parser extreu la llista de variables d'argument i utilitza la funció local \texttt{lambdifica} per niar dinàmicament constructors \texttt{Lam}. L'expressió es converteix internament en \texttt{Bind "suma" (Lam "x" (Lam "y" (...)))}. Aquest desugarejament d'arguments en cascada és la implementació directa de la transformació sintàctica formal formalitzada per Peyton Jones al capítol de traducció d'estructures algèbriques \cite{peyton1987implementation}.
    \item \textbf{Guàrdies Sintàctiques (\textit{Guards}):} Permet definicions basades en condicions alternatives mitjançant el caràcter \texttt{|}. El combinador absorbeix seqüencialment els fragments de text mitjançant \texttt{guardElement} i el cas de tancament obligatori \texttt{otherwise}. Llavors, utilitzant la funció \texttt{ififica}, transforma aquesta llista de parells \textit{(condició, resultat)} en un niu ordenat i seqüencial d'expressions \texttt{If} equivalents, seguint escrupolosament l'algorisme de traducció de guàrdies descrit a \cite{peyton1987implementation}.
\end{enumerate}

\section{Implementació de l'Arquitectura del Parser}

\subsection{Combinadors de Parsers Monàdics Primitius}
Un parser s'estructura conceptualment com una funció que rep una cadena de text (\texttt{String}) i en consumeix un prefix vàlid, retornant una llista de parells amb el valor semàntic extret i el fragment restant de la cadena sense analitzar:

\begin{lstlisting}[style=haskellstyle, caption=Definició monàdica del tipus Parser]
newtype Parser a = Parser {parse :: String -> [(a, String)]}
\end{lstlisting}

Per entendre com s'erigeix l'estructura del parser, és imprescindible analitzar els combinadors monàdics més bàsics de la implementació, els quals actuen com els maons primitius de l'edifici sintàctic:

\begin{lstlisting}[style=haskellstyle, caption=Combinadors monàdics primitius item i sat]
item :: Parser Char
item = Parser parseItem
  where
    parseItem [] = []
    parseItem (x : xs) = [(x, xs)]

sat :: (Char -> Bool) -> Parser Char
sat p = item >>= \x -> if p x then result x else zero
\end{lstlisting}

El combinador \texttt{item} és l'atòmic: consumeix incondicionalment el primer caràcter de la cadena. La veritable potència de la mònada es visualitza a \texttt{sat} (satisfar). Aquest combinador utilitza l'operador d'enllaç monàdic (\texttt{>>=}) per extraure el caràcter programat per \texttt{item} (anomenat \texttt{x}) i passar-lo a una funció de decisió en temps d'execució. Si el predicat \texttt{p x} és cert, retorna el caràcter encapsulat amb \texttt{result}; si falla, retorna \texttt{zero} (la llista buida, representant la fallida del \textit{backtracking}).
A partir d'aquesta composició monàdica elemental, es construeixen per abstracció combinadors funcionals d'alt nivell sense haver de tornar a manipular manualment l'estat de la cadena de text:

\begin{lstlisting}[style=haskellstyle, caption=Derivació de parsers de caràcters per abstracció]
digit :: Parser Char
digit = sat isDigit

lletra :: Parser Char
lletra = sat isLetter
\end{lstlisting}

Aquesta elegància declarativa s'estén a la instanciació de la classe \texttt{Alternative} (mitjançant l'operador \texttt{<|>}), que dota el sistema de la capacitat de tria prioritzada i recuperació d'errors. Els combinadors de repetició \texttt{mul} (zero o més vegades) i \texttt{mes} (una o més vegades) fan ús d'aquesta propietat per absorbir llistes de tokens recursivament.

\subsection{Gestió de Precedència i Jerarquia Sintàctica}
El repte de dissenyar una gramàtica per a expressions funcionals rau en l'ambigüitat inherent de les operacions infixes i les aplicacions juxtaposades. Per resoldre la \textbf{precedència d'operadors} sense recórrer a una taula de parsing externa o generadors complexos de tipus LR, s'ha seguit l'esquema de descens sintàctic jerarquitzat recomanat per Peyton Jones per a l'arquitectura de compiladors funcionalistes \cite{peyton1987implementation}.

L'arrel de l'anàlisi de qualsevol expressió és el combinador \texttt{expr}, el qual delega el control seguint un ordre estrictament creixent de precedència (de menys a més prioritat):
\begin{center}
\texttt{expr} $\rightarrow$ \texttt{lam} / \texttt{ifThenElse} $\rightarrow$ \texttt{operator} $\rightarrow$ \texttt{logicOr} $\rightarrow$ \texttt{logicAnd} $\rightarrow$ \texttt{comp} $\rightarrow$ \texttt{suma} $\rightarrow$ \texttt{term} $\rightarrow$ \texttt{ap} $\rightarrow$ \texttt{atom}
\end{center}

D'aquesta manera, l'operació multiplicativa \texttt{x * y} present en el cos de la nostra expressió guia és reconeguda pel combinador \texttt{term}, ramificant correctament l'AST abans que operacions de menor precedència com les sumes puguin capturar erròniament els operands, garantint l'arbre correcte descrit a \cite{peyton1987implementation}.

\subsection{Aplicació de Funcions i Currificació}
Com assenyala explícitament Peyton Jones \cite{peyton1987implementation}, l'aplicació de funcions en els llenguatges de la família del Càlcul Lambda és \textbf{associativa per l'esquerra}. L'expressió guia \texttt{(\textbackslash x y -> x * y) 2 3} no representa una funció rebent un vector de dos arguments, sinó que conceptualment s'agrupa obligatòriament com \texttt{((\textbackslash x y -> x * y) 2) 3}.

El combinador \texttt{ap} resol aquesta propietat geomètrica recollint primer un àtom base i acumulant en una llista un nombre indeterminat d'àtoms adjacents (\texttt{ys <- mul atom}). La funció interna \texttt{insert} aplica una reducció estructural per l'esquerra (\textit{left-fold}), anidant nodes constructors \texttt{App} de manera asimètrica:

\begin{lstlisting}[style=haskellstyle, caption=Associativitat per l'esquerra en el combinador d'aplicació]
ap :: Parser Expr
ap = do
    x <- atom
    ys <- mul atom
    return (insert (x:ys))
        where
            insert :: [Expr] -> Expr
            insert [e] = e
            insert exprs = (App (insert (init exprs)) (last exprs))
\end{lstlisting}

En contraposició, tal com dicta la sintaxi clàssica de Peyton Jones \cite{peyton1987implementation}, el combinador \texttt{lam} realitza l'operació inversa per a la definició de paràmetres: agrupa els identificadors de variables associant cap a la \textbf{dreta}. Així, la subexpressió \texttt{\textbackslash x y -> x * y} es desugareja automàticament durant el parsatge com a dues lambdes pures imbricades: \texttt{Lam "x" (Lam "y" (...))}.

\section{Arbre de Sintaxi Abstracta Resultant}
Una vegada el combinador final de la jerarquia ha consumit la totalitat de la cadena de caràcters de l'expressió guia, el parser retorna el corresponent Arbre de Sintaxi Abstracta (AST). 

Com es pot apreciar detalladament a la Figura \ref{fig:ast_or}, la semàntica de la currificació teòrica defensada a \cite{peyton1987implementation} esdevé perfectament visible en l'estructura de dades resultant: els nodes d'aplicació (\texttt{App}) s'apilen asimètricament cap a l'esquerra, col·locant l'argument terminal \texttt{Val 3} com a fill d'un node arrel que enllaça amb l'aplicació parcial prèvia. Aquesta fita clou la fase sintàctica i dota l'intèrpret de la representació formal requerida per al posterior mòdul d'inferència de tipus.

\begin{figure}[h]
    \centering
    \begin{forest}
      for tree={
        font=\ttfamily, 
        draw,           
        rounded corners,
        align=center,   
        edge={->},      
        s sep=10mm,     
        l sep=8mm       
      }
      [App
        [App
          [Lam
            ["x"]
            [Lam 
              ["y"]
              [App
                [App
                  [Var "*"]
                  [Var "x"]
                ]
                [Var "y"]
              ]
            ]
          ]
          [Val 2]
        ]
        [Val 3]
      ]
    \end{forest}
    \caption{Arbre de Sintaxi Abstracta (AST) generat pel parser per a l'expressió guia \texttt{(\textbackslash x y -> x * y) 2 3}.}
    \label{fig:ast_or}
\end{figure}